% TOP_PROBLEM: {"id": 2528, "title": "Kourovka Notebook Problem 21.19", "category_id": 17, "category": {"id": 17, "name": "group_theory", "display_name": "Group Theory", "description": "Problems about groups, group actions, representations, and related algebraic structures.", "slug": "group-theory", "order_index": 17, "created_at": "2026-07-31T15:26:25.671Z"}, "status": "open", "problem_number": "KOU-21.19", "set_id": 10}
% FIRST_POSTED: 2026-10-01
% ENTRY_KIND: statement_only
% UnsolvedMath ID 2528; source catalogue code KOU-21.19
% !TeX program = lualatex
% Definition and sources only; this file is not an LLM solution attempt.
\documentclass[11pt,a4paper]{article}
\usepackage[margin=25mm]{geometry}
\usepackage{fontspec}
\setmainfont{lmroman10-regular.otf}[BoldFont=lmroman10-bold.otf,
 ItalicFont=lmroman10-italic.otf,BoldItalicFont=lmroman10-bolditalic.otf]
\usepackage{amsmath,amssymb,mathtools}
\usepackage{unicode-math}
\setmathfont{latinmodern-math.otf}
\AtBeginDocument{\let\setminus\smallsetminus}
\usepackage{xurl,hyperref}
\hypersetup{colorlinks=true,urlcolor=blue}
\setcounter{secnumdepth}{0}
\setlength{\parindent}{0pt}
\setlength{\parskip}{5pt}
\setlength{\emergencystretch}{3em}
\begin{document}
\section*{Kourovka Notebook Problem 21.19}\label{2528}
{\small UnsolvedMath ID 2528 · KOU-21.19 · Group Theory · Kourovka Notebook - New Problems, Issue 21}
\subsection{Definitions and mathematical statement}
Work with definable groups and maps in a common sufficiently saturated
first-order structure; definability allows parameters. Morley rank is
the rank with $\operatorname{RM}(X)\ge0$ for nonempty definable $X$, and
$\operatorname{RM}(X)\ge k+1$ when $X$ contains infinitely many pairwise
disjoint definable subsets of rank at least $k$. A group has finite
Morley rank when the rank of its underlying definable set is finite.
It is connected if it has no proper definable subgroup of finite index.

Let $S$ be an infinite simple group of finite Morley rank, and suppose
every proper definable subgroup of $S$ is nilpotent. Simplicity means
that the only normal subgroups are $1$ and $S$. Nilpotence means that
the lower central series, obtained by repeatedly taking commutators
with the group, reaches $1$ after finitely many steps.
Let a nontrivial connected group $M$ of finite Morley rank act on $S$
by a definable map $M\times S\to S$, with each element inducing an
automorphism. Assume the action is faithful, so only $1\in M$ fixes
every element, and irreducible, so no definable subgroup
$1<H<S$ is invariant under all of $M$.

Must this action be equivalent to the conjugation action of $S$ on
itself? Explicitly, the required conclusion is that there are group
isomorphisms $\theta:M\to S$ and $\eta:S\to S$ intertwining the actions:
\[
 \eta(m\cdot s)=\theta(m)\eta(s)\theta(m)^{-1}
                       \quad(m\in M,\ s\in S).
\]
Thus identifying the acting group abstractly is insufficient without
identifying its action. The irreducibility hypothesis refers to all
definable subgroups, not only normal or connected ones.
\subsection{Short English statement}
Must a faithful irreducible definable action on an infinite simple finite-Morley-rank group with nilpotent proper definable subgroups be its conjugation action?
\subsection{Sources}
\begin{itemize}
\item[S1] ULAM.AI and the UnsolvedMath Contributors, supplied problems.json, record 2528 (KOU-21.19), Kourovka Notebook - New Problems, Issue 21. Catalogue identity and source transcription; CC BY 4.0.
\url{https://huggingface.co/datasets/ulamai/UnsolvedMath}.
\item[S2] The Kourovka Notebook, 21st edition, new problems, Problem 21.19. Expanded formulation of the supplied catalogue statement.
\url{https://arxiv.org/pdf/1401.0300v47}.
\end{itemize}
\end{document}
